\section{Discussion}
\label{sec:discussion}

\subsection{What the Results Show}
The per-class penalty is the part of the reward that decides which tasks are protected. Proposition 1 shows why. A learner maximizes the expected reward, which is the same as minimizing the sum of $(w_c + \beta)$ times the misses of each class. With one penalty for all classes, this is the overall miss rate, and the DP policy of the flat reward is 65\% to 118\% above the optimum in every system, however well it is learned. With the operator weights as penalties, the reward is close to the operator's objective, and with the corrected penalties of Eq.~\eqref{eq:corrected} it is the operator's objective. Trained PPO agents follow the same order in every system, and DQN agents in the nominal system.

The exact benchmark also shows how much the learners leave on the table. In the nominal system the best agent is 10\% above the optimum and beats a tuned rule. In the larger systems the agents are 36\% to 67\% above the optimum. They beat the simple reservation rule in two of the three systems, but the node-set rule, which knows the latency model and is tuned on the exact model, is better in all three. Without an optimal reference, an agent that beats join-shortest-queue, an SLA rule and a class-blind version of itself would look strong in all four systems. The gap to the optimum would stay invisible.

The rebuilt published design points to the same cause from another side. Its adaptive tiers raise every penalty whose target is missed. When the targets cannot all be met, the penalties drift toward each other, and the agent drifts toward the class-blind policy. In the nominal system this left it 41\% to 60\% above the optimum, between the fixed weighted reward and the flat reward. Setting penalty targets is safe only when it is known which targets are reachable, which the exact frontier shows.

\subsection{When to Use Which Policy}
If the latency model is known and the system is small enough for the exact model, the DP policy is the best choice, because it is optimal by construction. A tuned reservation rule is a good second choice. It needs only the latency model and a search over a few thresholds, and under every change of the nominal system that was tested it stayed between 6\% and 16\% above the new optimum without being tuned again. A learned agent is useful when the latency model is not known or changes, or when the system is too large for the exact model. Table~\ref{tab:scaling} shows that this happens quickly, since every extra edge node multiplies the stored transitions of the exact model by 14 to 24. In that case the results of this paper give three practical rules. Use a per-class penalty, because a class-blind penalty optimizes the wrong objective. Fold the bonus for a met limit into the penalties with Eq.~\eqref{eq:corrected}, so that the reward has the operator's weights. Tune and train long enough, and check the result against a strong rule.

\subsection{Threats to Validity}
The systems are synthetic. Latencies follow a linear queue model, the cloud is described only by the chance to meet each limit, and there is no energy or transmission model. Real traces were not available to the authors. The results show what the reward does in systems where the answer can be checked exactly, not how large the effect is in a deployed system. The SLA cost uses the operator weights, which the weighted reward also uses. The per-class miss rates, the trade-off curve of Fig.~\ref{fig:frontier} and the weight sweep of Table~\ref{tab:ratios} let a reader with other weights judge every policy. The agents use standard multilayer networks from one library, tuned with six configurations in the nominal system and three in each larger system. Other architectures, longer training or larger searches could bring them closer to the optimum. The rebuilt published design follows the equations of its paper, but parameters the paper does not give had to be chosen, and two step sizes were tested for that reason. All exact values rely on the exact model, which the test suite checks against the simulator.

\section{Conclusion}
\label{sec:conclusion}

This paper studied the reward of deep RL task offloading in the IoT edge-cloud continuum against an exact optimum. Four systems with up to four heterogeneous edge nodes and three task classes were solved by DP, which gave the optimal policy, the optimal trade-off between the classes and an exact score for every trained agent. A reward with a bonus for a met limit ranks policies by the per-class penalty plus the bonus. A class-blind penalty therefore minimizes the overall miss rate and stays 65\% to 118\% above the optimum with a perfect learner, while a corrected per-class penalty reaches the optimum. Trained PPO agents show the same effect in every system and DQN agents in the nominal system, and in the nominal system tuned PPO comes within 10\% of the optimum and beats a tuned reservation rule in 9 of 10 seeds. In the larger systems the agents remain 36\% to 67\% above the optimum, and a node-set rule tuned on the exact model is better than every agent. A rebuilt published design with adaptive penalty tiers stays between the weighted and the class-blind reward, because its targets cannot all be met and its penalties drift together. Future work will use the exact benchmark to test learners that use the structure of the system, approximate DP for systems beyond five edge nodes, and models of energy and transmission as in \cite{hu2021,zhang2024}.
